\scanpage{42}
\unsure{tel|tels} que
\[
\left\{\begin{aligned}\text{a) }&g(e_B)=1_{C_n}\\
\text{b) }&g(B)\text{ engendre }C_n.
\end{aligned}\right. \tag*{(67)}
\]
Ceci dit, considérons \deleted{l'application \unsure{fonction}} \inserted{le morphisme}
\[
F_n:\check{\mathbb V}(C_n)\longrightarrow\check{\mathbb V}(\mathcal O_S) \tag*{(68)}
\]
défini par
\[
 F_n(x)=\unsure{\mathcal N|\mathscr N|N}_{C_n/\mathcal O_S}(x)\;(=\mathcal N_{\Omega/S}(x)), \tag*{(69)}
\]
c'est une $n$-forme, et on peut considérer
sa ``restriction'' $F_n\circ g$. Ceci dit, on
aura (sauf erreur)
\[
 f_n=\deleted{\unsure{(-1)^n}}F_n\circ g\ \inserted{=\deleted{\unsure{(-1)^n}}\mathcal N_{C_n/\mathcal O_S}\circ g}\quad\text{dans }\Gamma\underline{\mathrm{Sym}}^n(\check B=V) \tag*{(70)}
\]
[ce qui montre que la correspondance
entre $n$-diviseurs sur $F$, \deleted{et} i.e. syst. $(C_n,g)$
satisfaisant (67), et $n$-formes définie par
(70), est bijective].

Pour vérifier (70), OPS que $B$ est muni d'
une section $t$ telle que $(e_B,\unsure{T|t})$ forment une
base. Alors $C_n$ sera engendré sur $\mathcal O_S$ par \inserted{l'image $t$ de} $T$,
donc aura la base $1,t,\ldots,t^{n-1}$, et sera défini
par
\[
 C_n\simeq\mathcal O_S[T]/(f)\mathcal O_S[T] \tag*{(71)}
\]
où $f$ est un polynôme unitaire \inserted{en $T$} de \uline{degré} $n$,
\[
 f(T)=\deleted{\unsure{a_0+}}T^n+a_1T^{n-1}+a_2T^{n-2}\,\text{---}\,+a_n
\]
\deleted{qui n'est autre que $f_n(t)$}
On \deleted{\unsure{a}} vérifie de suite que l'on a $N$
\[
\begin{aligned}
 M_n&=\mathcal O_S\,g_n\qquad\deleted{\unsure{g_n=e_B^n+a_1 e_B^{n-1}T+\cdots+a_n\in\Gamma\mathrm{Sym}^n B|q_n=e_B^n+a_1e_B^{n-1}T+\cdots+a_n\in\Gamma\mathrm{Sym}^n B}}\\
 \deleted{(72)}\quad&\deleted{\unsure{f_n(u e_B+vT)=u^n+a_1u^{n-1}v+a_2u^{n-2}v^2+\cdots+a_nv^n}}\\
 &\unsure{B}\quad g_n=\deleted{\unsure{a_n}}T^n+a_{n-1}e_BT^{n-1}+\cdots+a_ne_B^n\in\Gamma\mathrm{Sym}^n B
\end{aligned}\tag*{(72)}
\]
% archive sheet number: 41
